\chapter{Vaginismus}

\section*{Definition}
Vaginismus is a persistent or recurrent difficulty with vaginal entry during penetrative attempts (sexual intercourse, tampon insertion, speculum exam) due to involuntary pelvic floor muscle contraction and/or fear of pain or penetration. It is classified under genito-pelvic pain/penetration disorder in DSM-5 when present for at least 6 months.

\section*{Pathophysiology}
Vaginismus involves a conditioned reflex arc where anticipation of pain triggers involuntary spasm of the pubococcygeus and other pelvic floor muscles. This is mediated by hyperexcitability of the sacral spinal reflex and reduced descending inhibitory control from the prefrontal cortex and anterior cingulate, reinforced by anticipatory anxiety and avoidance.

\section*{Etiology}
\begin{itemize}
  \item \textbf{Primary (lifelong):} No history of successful penetration
  \item \textbf{Secondary (acquired):} Develops after a period of normal function
  \item \textbf{Sexual trauma:} Sexual abuse, assault, or rape (most common risk factor)
  \item \textbf{Medical factors:} Vulvodynia, vestibulodynia, vaginitis, childbirth trauma (episiotomy, perineal tear), pelvic surgery, radiation, endometriosis, lichen sclerosus, postmenopausal atrophy
  \item \textbf{Psychological:} Negative sexual attitudes, strict religious upbringing, anxiety disorders, fear of pain, catastrophizing, relationship conflict
  \item \textbf{Other:} Vaginismus can occur without identifiable organic cause (idiopathic)
\end{itemize}

\section*{Indications for Evaluation}
Seek care if:
\begin{itemize}
  \item Difficulty with vaginal penetration causes distress or prevents sexual activity
  \item The individual desires to use tampons or undergo gynecological examination but cannot
  \item Associated with infertility, relationship problems, or avoidance of all gynecologic care
\end{itemize}

\section*{Related Symptoms}
\begin{itemize}
  \item Dyspareunia, vulvodynia, pelvic floor hypertonicity, pelvic pain, anxiety, avoidance of intimacy, sexual dysfunction
\end{itemize}
\textbf{Note:} Vaginismus is a treatable condition; pelvic floor physical therapy combined with psychological treatment yields high success rates.

\section*{Self-Care / Home Management}
\begin{itemize}
  \item Use graduated vaginal dilators (starting with the smallest size) with water-based lubricant
  \item Practice Kegel exercises followed by active relaxation (Kegel squeeze--release)
  \item Deep breathing and progressive muscle relaxation during penetration attempts
  \item Do not force penetration; stop and relax when the spasm occurs
\end{itemize}

\section*{Diagnostic Approach}
\begin{itemize}
  \item Clinical history: onset, precipitating factors, pain characteristics, prior trauma
  \item Gynecologic exam: if the patient can tolerate, assess pelvic floor muscle spasm and tenderness; often deferred to after treatment
  \item Assess for secondary causes: wet mount, culture, biopsy if lesions present
  \item Psychological assessment for trauma history, anxiety, and relationship factors
\end{itemize}

\section*{Treatment / Management Overview}
\begin{itemize}
  \item First-line: Graduated vaginal dilator therapy combined with pelvic floor physical therapy
  \item Botulinum toxin A injection into the pubococcygeus muscle for severe, refractory cases
  \item Topical lidocaine or compounded baclofen/diazepam suppositories
  \item Cognitive-behavioral therapy: exposure therapy, cognitive restructuring of catastrophic beliefs
  \item Sex therapy with partner involvement to address relationship dynamics
  \item Psychological treatment for trauma-focused therapy if sexual abuse history is present
\end{itemize}

\clearpage
